%==============================================================================
%  Multimedia Appendix 2 — Prompt templates
%  Companion to main_rt.tex. Compiles standalone; export to DOCX for submission.
%  INTEGRITY: templates below are the deployed-system templates in normalized
%  form. Replace \PH{...} slots with the exact production strings before
%  submission so the appendix is byte-faithful to what was run.
%==============================================================================
\documentclass[10pt]{article}
\usepackage[utf8]{inputenc}
\usepackage[T1]{fontenc}
\usepackage[margin=1in]{geometry}
\usepackage{times}
\usepackage{booktabs}
\usepackage{array}
\usepackage{xcolor}
\usepackage{caption}
\usepackage{fancyvrb}
\usepackage{enumitem}
\usepackage[colorlinks=true,allcolors=blue!55!black]{hyperref}

\newcommand{\PH}[1]{{\color{red}\textbf{[#1]}}}
\newcolumntype{L}[1]{>{\raggedright\arraybackslash}p{#1}}
\captionsetup{labelsep=period,font=small,justification=justified,
              singlelinecheck=false,skip=4pt}
\setcounter{secnumdepth}{0}
\pagestyle{plain}

\begin{document}

\begin{center}
{\large\bfseries Multimedia Appendix 2. Prompt templates}\\[2pt]
{\footnotesize Companion to: \emph{Large Language Models for Automating Hospital
Statistical Criteria Conversion to Clinical Data Warehouse Queries: Validation
and Evaluation Study}}
\end{center}

\vspace{6pt}
\noindent
This appendix lists every prompt used by the automated pipeline. Prompts are
grouped into (1) \textbf{pipeline task prompts}, one per module, used in all
experiments, and (2) \textbf{the five prompting strategies} that were varied in
the factorial hallucination experiment. All prompts were issued with
deterministic decoding (temperature 0, top-$p$ 1). Braces \verb|{...}| denote
runtime substitution slots.

%==============================================================================
\section{Part 1 --- Pipeline task prompts}
%==============================================================================

\subsection{P1. Segmentation (preprocessing stage 1)}
\begin{Verbatim}[fontsize=\footnotesize,frame=single,framesep=4pt]
System: You split a free-text hospital statistical request into atomic
statistical criteria. Preserve Boolean connectives (AND/OR/NOT) and any
grouping hierarchy (by month, by department, by ward). Do not resolve,
interpret, or normalize anything at this stage. Do not drop text.
Output: a JSON array; each element is
  {"id": <int>, "text": <verbatim span>, "connective": "AND"|"OR"|"NOT"|null,
   "parent_id": <int|null>}
Request:
{raw_request_text}
\end{Verbatim}

\subsection{P2. Filtering (preprocessing stage 2)}
\begin{Verbatim}[fontsize=\footnotesize,frame=single,framesep=4pt]
System: You mark which atomic criteria are queryable against a relational
clinical data warehouse. Mark as NON-QUERYABLE: formatting and delivery
instructions, deadlines, narrative justification, requester contact details,
and any statement that expresses a preference rather than a data condition.
Do not delete; only label, with a one-clause reason.
Output: JSON array of {"id": <int>, "queryable": true|false, "reason": <string>}
Criteria:
{segmented_criteria_json}
\end{Verbatim}

\subsection{P3. Simplification (preprocessing stage 3)}
\begin{Verbatim}[fontsize=\footnotesize,frame=single,framesep=4pt]
System: You rewrite each queryable criterion into one short, unambiguous
sentence. Normalize temporal expressions to explicit closed intervals in the
form [YYYY-MM-DD, YYYY-MM-DD). Resolve relative dates against {report_date}.
Keep every clinical or administrative term verbatim -- never substitute a
synonym. Never add a condition that is not present in the input.
Output: JSON array of {"id": <int>, "simplified": <string>}
Criteria:
{filtered_criteria_json}
\end{Verbatim}

\subsection{P4. Information extraction (seven structured elements)}
\begin{Verbatim}[fontsize=\footnotesize,frame=single,framesep=4pt]
System: From each simplified criterion, extract exactly seven element types:
  1 Business Terms      -- institutional terms as written
  2 Metadata Systems    -- which knowledge source should resolve the term
                           (warehouse schema | metric-criterion dictionary |
                            historical SQL | exception rules)
  3 Field Identifiers   -- only if the criterion names a column explicitly
  4 Values              -- literal comparison values, with unit if present
  5 Attributes          -- aggregation, grouping, ordering, distinctness
  6 Temporal            -- the counting time criterion, if stated
  7 Negation            -- scope of any exclusion
Leave a slot null rather than inferring it. Do NOT emit column names here.
Output: JSON array keyed by criterion id.
Criteria:
{simplified_criteria_json}
\end{Verbatim}

\subsection{P5. Field and criterion mapping (few-shot)}
\begin{Verbatim}[fontsize=\footnotesize,frame=single,framesep=4pt]
System: You map each business term to EXACTLY ONE warehouse field and, when the
term denotes an indicator, to EXACTLY ONE counting criterion. You may only
choose from the retrieved candidates supplied below. You may not invent a table
name, a column name, or a criterion.

Hard rules:
  - If two or more candidate criteria are defensible, do NOT choose. Emit a
    TO-CONFIRM item listing the candidates and what distinguishes them.
  - If no candidate is defensible, emit TO-CONFIRM with reason "no candidate".
  - Timestamps are never interchangeable. Room-in time, surgery-completion
    time, discharge time and application time denote different counting
    populations.
  - Report the layer (DL/DC/MDR) of the field you chose.

Retrieved candidates (top-k=8 per term, hybrid BM25 + dense):
{retrieved_candidates_json}

Worked examples:
{few_shot_examples}

Terms to map:
{extracted_terms_json}

Output: JSON array of
  {"term": <string>, "table": <string>, "column": <string>, "layer": <string>,
   "criterion_id": <string|null>, "confidence": <0-1>,
   "to_confirm": <null | {"candidates": [...], "reason": <string>}>}
\end{Verbatim}

\subsection{P6. SQL generation (few-shot, sub-structure decomposition)}
\begin{Verbatim}[fontsize=\footnotesize,frame=single,framesep=4pt]
System: You generate Presto SQL for hospital statistical reporting. Use ONLY
the tables, columns and value mappings in the resolved mapping below; never
invent identifiers.

Mandatory in every generated statement:
  - every joined table carries its logical-deletion filter isdeleted = '0'
  - every multi-campus table carries its organization-code filter
  - all warehouse columns are VARCHAR: cast explicitly before any date or
    numeric comparison; never rely on implicit coercion
  - quote non-ASCII column names with double quotes

Emit the query one sub-structure at a time, in this order, each with a
one-line justification citing the mapping entry it comes from:
  SELECT / FROM-JOIN / WHERE / GROUP BY / ORDER BY

Then emit:
  1) the assembled statement
  2) field provenance: for every selected column, its source table, layer,
     and the mapping entry id
  3) any TO-CONFIRM items carried forward unresolved

Resolved mapping:
{mapping_json}
Historical reference queries for this indicator family:
{historical_sql_examples}
Request:
{simplified_criteria_json}
\end{Verbatim}

%==============================================================================
\section{Part 2 --- The five prompting strategies (factorial experiment)}
%==============================================================================

\noindent
In the factorial experiment (24 tasks $\times$ 8 models $\times$ 5 strategies),
prompt P6 above was replaced by one of the following five variants. Everything
else in the pipeline was held constant.

\begin{table}[h]\centering
\caption{The five prompting strategies varied in the factorial experiment.}
\footnotesize
\begin{tabular}{@{}L{3.0cm}L{11.5cm}@{}}
\toprule
Strategy & Instruction body appended to the task prompt\\
\midrule
zero\_shot &
\emph{(no additional instruction)} --- the resolved mapping and the request are
supplied and the model is asked directly for the Presto statement.\\
\addlinespace
structured\_ approach &
``Work through the query one clause at a time. First decide the FROM-JOIN path,
then the WHERE conditions, then the SELECT list, then GROUP BY, then ORDER BY.
State your reasoning for each clause before writing it, then assemble.''\\
\addlinespace
explicit\_ uncertainty &
``If any field or counting criterion is ambiguous, do not guess. Emit
`TO-CONFIRM: <question>' in place of that clause and continue. A query with an
honest TO-CONFIRM is preferred over a complete query that may count the wrong
population.''\\
\addlinespace
validation\_ focused &
``Before returning, verify each of: every identifier appears in the supplied
mapping; every joined table has its logical-deletion filter; every multi-campus
table has its organization-code filter; every date comparison has an explicit
cast. Report the checklist result, then the query.''\\
\addlinespace
error\_aware &
``The target instance has partial metadata coverage: many production columns are
absent, and some flag columns are entirely unpopulated. Common failures are (a)
selecting a structurally valid but unpopulated flag, (b) substituting a
different timestamp for the requested counting criterion, and (c) leaving a
placeholder literal. Avoid all three; if unavoidable, say which one you hit.''\\
\bottomrule
\end{tabular}
\end{table}

%==============================================================================
\section{Part 3 --- Reproduction notes}
%==============================================================================
\begin{itemize}[leftmargin=1.4em,itemsep=1pt,topsep=2pt]
  \item Decoding: temperature 0, top-$p$ 1, for all 8 models, all 5 strategies.
  \item Retrieval: hybrid BM25 + dense sentence embeddings, top-$k$=8 per term,
        over the metadata knowledge base; identical for every model.
  \item Few-shot exemplars for P5/P6 were drawn only from the 30-task
        development set and never from the 24-task evaluation set.
  \item Exact model builds, endpoint versions, and serving configuration:
        \PH{fill from the run manifest on the compute cluster}.
  \item Random seeds are not applicable under deterministic decoding; the
        SynHDW generator seed is \PH{fill}.
\end{itemize}

\end{document}
